\documentclass[10pt,aspectratio=169]{beamer}
\usepackage[portuguese]{babel}
\usepackage{fontspec,unicode-math}
\setmainfont{Noto Sans}
\setsansfont{Noto Sans}
\setmathfont{STIX Math}
\usepackage{amsmath}
\usepackage{xcolor}
\definecolor{InsperRed}{HTML}{a01924}
\definecolor{Ink}{HTML}{232529}
\setbeamercolor{frametitle}{fg=Ink}
\setbeamercolor{structure}{fg=InsperRed}
\setbeamercolor{normal text}{fg=Ink,bg=white}
\setbeamerfont{frametitle}{size=\large,series=\bfseries}
\setbeamerfont{title}{size=\LARGE,series=\bfseries}
\setbeamertemplate{navigation symbols}{}
\setbeamertemplate{itemize item}{\color{InsperRed}$\bullet$}
\setbeamertemplate{footline}{%
  \hbox{\begin{beamercolorbox}[wd=.84\paperwidth,ht=2.7ex,dp=1.2ex,leftskip=.7cm]{normal text}\scriptsize Inteligência Artificial para Decisão · Insper · 2027-1\end{beamercolorbox}%
  \begin{beamercolorbox}[wd=.16\paperwidth,ht=2.7ex,dp=1.2ex,rightskip=.7cm plus1fil]{normal text}\hfill\scriptsize\insertframenumber\, /\, \inserttotalframenumber\end{beamercolorbox}}}
\setbeamersize{text margin left=12mm,text margin right=12mm}
\begin{document}
\begin{frame}[t]{}
\vspace{14mm}\textcolor{InsperRed}{\bfseries AULA 06}\par\vspace{4mm}{\LARGE\bfseries Composição de aplicações com LLMs}\par\vspace{6mm}{\large Marketing Analítico}\par\vspace{12mm}{\color{gray}Paulo C. Marques F. · 2027-1}
\end{frame}

\begin{frame}[t]{Objetivos e verificação}
\small
\begin{itemize}
\item Projetar fluxos compostos, identificar dependências entre ferramentas e distinguir informação observada, cálculo derivado e inferência do modelo.
\item Resolver manualmente os exemplos do capítulo, identificando hipóteses e consequências.
\item Especificar a aplicação no Codex gráfico e confrontar ao menos um resultado com evidências.
\end{itemize}
\end{frame}

\begin{frame}[t]{Organização do encontro}
\small
\begin{itemize}
\item Apresentação dos fundamentos e resolução de problemas quantitativos no papel.
\item Laboratório supervisionado e testes produzidos por prompts na interface gráfica do Codex, sem terminal.
\item Atividade individual com dados diferentes e depósito no GitHub antes do quiz seguinte.
\end{itemize}
\end{frame}

\begin{frame}[t]{Aplicações como composição de operações}
\small
\begin{itemize}
\item Uma aplicação composta transforma a saída de uma operação em entrada de outra.
\item Se prepara avaliações, recupera textos semelhantes e combina indicadores comerciais, o resultado é , em que reúne preços e participações provenientes de uma base independente.
\item A composição precisa respeitar dependências, tipos, unidades e proveniência; uma operação isolada correta não garante que o sistema completo esteja correto.
\item Um grafo acíclico de tarefas torna explícitas as relações entre etapas.
\end{itemize}
\end{frame}

\begin{frame}[t]{Aplicações como composição de operações — continuação}
\small
\begin{itemize}
\item Para cada aresta, o aluno deve perguntar: ``Qual dado passou?
\item Campos com mesma aparência mas sem mesmo significado não podem ser combinados automaticamente.
\end{itemize}
\end{frame}

\begin{frame}[t]{Proveniência e temporalidade}
\small
\begin{itemize}
\item Um relatório sobre posicionamento pode reunir avaliações coletadas em determinado mês, preços vigentes em outra data e participação de mercado aferida em um período anual.
\item Somar essas informações sem registrar suas datas cria uma aparência de precisão que não existe.
\item Toda saída deve identificar de onde veio cada número, qual transformação o produziu e quais premissas são necessárias para interpretá-lo.
\item Considere preços de R\textbackslash\{\} 150.
\end{itemize}
\end{frame}

\begin{frame}[t]{Proveniência e temporalidade — continuação}
\small
\begin{itemize}
\item A diferença relativa a 150 é ; relativa a 120 é .
\item Um assistente que afirma ``a marca A é 25\% mais barata que B'' escolheu a base errada se pretende medir desconto em relação a B.
\item A noção de porcentagem exige denominador explícito.
\end{itemize}
\end{frame}

\begin{frame}[t]{Falhas de integração}
\small
\begin{itemize}
\item Quatro erros merecem especial atenção: (i) ferramenta não executada; (ii) argumento incorreto ou de unidade incompatível; (iii) retorno correto atribuído à entidade errada; (iv) inferência não sustentada apresentada como…
\item Um rastro de execução útil permite localizar essas falhas.
\item A participação de mercado não pode ser estimada somente a partir da proporção de avaliações de determinada marca: o número de reviews depende de cobertura, propensão a avaliar e exposição.
\item Comparar a percepção textual com participação já observada é legítimo; deduzir participação diretamente do volume de reviews, sem desenho amostral apropriado, não é.
\end{itemize}
\end{frame}

\begin{frame}[t]{Verificação de ponta a ponta}
\small
\begin{itemize}
\item O laboratório deve produzir uma tabela final acompanhada de um registro de origens.
\item Testes unitários verificam cada ferramenta; testes de integração verificam a correspondência correta entre marca, data e indicador.
\item O aluno deverá fornecer pelo menos um exemplo em que cada componente funciona isoladamente, mas a associação final é incorreta.
\end{itemize}
\end{frame}

\begin{frame}[t]{Exercício para resolver no papel 1}
\small
\small\textbf{6.1 O desconto aparente}\par\vspace{4mm}A loja A oferece um produto a R\textbackslash\{\} 150 na mesma data. Calcule a diferença percentual tomando A como referência e tomando B como referência. Por que divulgar uma única porcentagem sem declarar a base pode enganar?
\end{frame}

\begin{frame}[t]{Exercício para resolver no papel 2}
\small
\small\textbf{6.2 Menção não é market share}\par\vspace{4mm}Em 80 avaliações, 48 mencionam `durabilidade''. Uma apresentação afirma: `A marca tem 60\% do mercado de durabilidade''. Calcule a proporção observada e explique precisamente a inferência inválida.
\end{frame}

\begin{frame}[t]{Laboratório e caso aplicado}
\small
\begin{itemize}
\item Área de aplicação: Marketing Analítico.
\item Mapear em papel quais conclusões cada fonte permite e quais não permite.
\item Pedir ao Codex painel que una atributos mencionados, preços datados e participações somente quando o levantamento fornecer o denominador.
\item Forçar um caso de dados conflitantes por data, unidade ou marca e verificar se o painel bloqueia comparação indevida.
\end{itemize}
\end{frame}

\begin{frame}[t]{Prompt de desenvolvimento — exemplo}
\small
\small\textbf{Instrução fornecida ao Codex:}\par\vspace{5mm}Amplie o painel de Marketing para combinar embeddings de avaliações com preços e dados explícitos de pesquisa de mercado. Exiba origem, data, denominadores e unidade de cada indicador; impeça inferir market share apenas de avaliações.
\end{frame}

\begin{frame}[t]{Critérios de aceitação e validação}
\small
\begin{itemize}
\item Toda afirmação numérica tem fonte; comparações usam datas/unidades compatíveis; participação de mercado nunca é inferida de frequência de comentários; conflito gera alerta visível.
\item Inspecionar argumentos, saídas e rastros das ferramentas quando houver chamadas externas.
\item Comparar uma saída com a resolução manual do exemplo conceitual.
\end{itemize}
\end{frame}

\begin{frame}[t]{Atividade individual em casa}
\small
\begin{itemize}
\item Conjunto inteiramente novo: avaliação de duas marcas fictícias de aspiradores, preços em duas datas e levantamento com 60 respostas de pesquisa sobre atributos.
\item Resolver as questões de interpretação, com cálculos e evidência documental.
\item Depositar no GitHub os arquivos solicitados; a entrega habilita o quiz, sem valer nota.
\end{itemize}
\end{frame}

\begin{frame}[t]{Preparação para a avaliação escrita}
\small
\begin{itemize}
\item Exercícios do caderno: 6.1 O desconto aparente; 6.2 Menção não é market share; 6.3 Uma evidência com duas datas.
\item Formalizar o raciocínio e interpretar resultados; não será exigida escrita manual de código.
\item No quiz seguinte serão avaliados conceitos desta aula e da atividade feita em casa.
\end{itemize}
\end{frame}

\end{document}
